\section{Related work}
\label{sec:related-work}

\paragraph{Static factorization and coordinate dependence.}
Under $X=\Theta^\top$, $W=B^\top$, and $H=A^\top$, our static model is
unconstrained SSMF. The uniformizing path of
\citet[Sec.~IV-A]{vuthanh2023bounded} is exactly the one above after
$s=\delta/(1+\delta)$; positive static recovery requires additional
geometric conditions~\citep{fu2018identifiability,abdolali2021simplex}.
Low-rank factor gauges, ordinary-update dependence on coordinates, and
gauge-sensitive merging are also established
\citep{yen2025lora,zheng2026loras,gong2026twistedmerge}.
Our auxiliary ASLoRA experiment applies this perspective to its documented
first-merge rule~\citep{hu2024aslorav2}.

\paragraph{Optimizer response and identifiability.}
Jacobian--Gram operators and induced factor flows are standard
\citep{jacot2018neural,marcotte2026intrinsic}.
\citet[Lem.~D.2]{varre2026gradient} already derive the single-router
squared-covariance softmax flow; Euclidean factor isometries and compatible
softmax permutation/shift updates are studied by
\citet{singh2026loss,lau2026symmetry}.
Function-level and conditional-density softmax-MoE identifiability and
inverse bounds use different observations
\citep{tran2025linear,nguyen2023demystifying}.
Our question is a converse at two fixed factorizations of the same finite
product: does equal complete response exclude every non-permutation gauge?
It is metric-specific; natural-gradient or quotient geometries can instead
remove chart dependence~\citep{vanoostrum2023invariance,mishra2014fixedrank}.

\paragraph{The tools behind the inverse.}
Orthogonal-to-permutation mechanisms, joint-diagonalizer uniqueness and
sensitivity, and robust tensor uniqueness precede our final algebraic steps
\citep{dorrell2026convex,afsari2008sensitivity,bhaskara2014uniqueness}.
In particular, \citet[Thm.~5]{bhaskara2014uniqueness} already give a linear
robust inverse when their conditioning parameters are fixed. That work establishes linear dependence and global robust uniqueness.
Our quantitative analysis tracks the stable reduction from response observations to diagonal data,
eliminates scaling through the simplex constraints, and controls both
original factors with explicit constants. Matrix--vector and structured
operator probing are likewise established
\citep{novak2022fast,chiu2012matrix,halikias2024structured,amsel2026query}.
The noisy-product result combines a uniform bound outside the observed
subspace with this response inverse. Whether $K$ queries are optimal remains open.
Appendix~\ref{app:literature} maps the closest results, the wider sharing
literature, and unresolved full-text access limitations.


\paragraph{Clustering and preserving training dynamics.}
Kernel $k$-means already expresses within-cluster distortion as a
Gram-matrix trace objective~\citep{dhillon2004kernel}.
Neural Tangent Transfer matches both initial outputs and empirical training
kernels to construct trainable sparse networks~\citep{pmlr-v119-liu20o}.
These works establish Gram-based clustering, kernel preservation, and
the tradeoff between matching initial states and training dynamics. Here
we reduce the \emph{complete softmax-factor response} exactly to balanced
router clustering: its local covariance blocks contribute a partition-independent
term. The product-preserving decision counterexample and the explicit
transverse crossing delimit the information and dynamical consequences for
this shared-basis model. A standard Gronwall bound relates the two distortion objectives along
training trajectories.
